\begin{table}[!htb]
\centering\small
\caption{Closest work and the precise difference.  The rightmost column states
what we do \emph{not} claim, so that standard components are not restated as
novelty.}
\label{tab:literature}
\begin{tabular}{p{3.1cm}p{3.6cm}p{4.4cm}p{2.4cm}}
\toprule
Work & What it already gives & Our relation & Not claimed \\
\midrule
Huang--Zhu (arXiv:2102.05469) & purchases observation opportunities in an LQG pursuit--evasion game & we budget \emph{command revisions} on a public calendar; information is observed continuously, action is not & neither model subsumes the other \\
Hogeboom-Burr--Y\"uksel (arXiv:2005.06673) & more information cannot hurt the informed player (zero-sum comparison) & our inclusion lemma is a finite special case; used as a correctness condition & not claimed as new \\
Chen--Waggoner (arXiv:1703.08636) & substitutes vs complements of signals depend on the decision problem & we give a binary-action counterexample showing calendars need not be submodular & extends the caution to revision calendars \\
McMahan et al. (ICML 2003) & robust planning against adversarial costs via response oracles & our pricing recurrence specialises one-sided generation to sealed-block histories & specialisation, not a new oracle theory \\
Bo\v{s}ansk\'y et al. (JAIR 2014) & exact double oracle for extensive-form zero-sum games & we optimise the information structure (calendar), not a strategy inside a fixed one & no raw node-count comparison across different games \\
McAleer et al. (XDO, 2021) & double oracle with mixing at information states & our root-policy generation is not XDO and carries no polynomial iteration guarantee & oracle tailored to a small history lattice \\
Kroer--Sandholm (NeurIPS 2018) & unified abstraction bounds for extensive-form games & our deletion certificate bounds loss of a specified feasible calendar policy by an exact full response & narrower than their general theorem \\
\bottomrule
\end{tabular}
\end{table}
